\section{Discussion}
\label{sec:discussion}

\subsection{The Floor of Vendor-Neutral LLM Development}
\label{sec:discussion:floor}

MiniMind demonstrates that the \emph{algorithmic} layer of LLM development
is largely vendor-neutral today.  Every training stage — from BPE tokenizer
through agent RL — runs on both H100 and MI350 without changes to the
trainer code, using only the portability layer in \texttt{trainer/device\_utils.py}
(\texttt{DeviceCtx}, \texttt{dist\_backend()}, \texttt{default\_device\_str()}).
PyTorch-ROCm's aliasing of the \texttt{torch.cuda.*} surface is sufficiently
mature that the remaining divergences are confined to:

\begin{enumerate}
\item \textbf{SDPA backend selection.} \texttt{F.scaled\_dot\_product\_attention}
  dispatches to cuDNN FlashAttention (CUDA) or to hipBLASLt-fused attention /
  Composable Kernel (ROCm) depending on the runtime.  The same PyTorch API
  surface is used in both cases; the performance and numerical properties of
  the underlying kernels differ.
\item \textbf{\texttt{torch.compile} warm-up time.}  First-run Triton-ROCm
  kernel compilation is slower than CUDA due to the additional LLVM→HSACO
  step.  Subsequent runs are cached.
\item \textbf{FP8 availability.}  NVIDIA exposes FP8 via Transformer Engine
  with a stable API; AMD's FP8 path via hipBLASLt is available but not yet
  wrapped in a \texttt{torch.nn.Linear}-compatible interface in stock PyTorch.
\item \textbf{Inference ecosystem maturity.}  vLLM and SGLang maintain
  ROCm branches, but the feature gap (e.g.\ speculative decoding, RadixAttention
  prefix caching) between the CUDA-primary and ROCm paths closes with each
  release cycle.
\end{enumerate}

\subsection{Where Divergence Remains}
\label{sec:discussion:divergence}

\paragraph{Custom kernels.}
Projects that rely on hand-written CUDA kernels (Flash-Attention~3, DeepSpeed
Triton extensions, Liger-Kernel) face a porting cost on ROCm.  MiniMind
deliberately avoids this by relying exclusively on \texttt{F.scaled\_dot\_product\_attention}
and standard PyTorch operations.  This is a deliberate design choice that
trades potential peak performance for portability.

\paragraph{MoE token dispatch.}
Expert routing involves a scatter/gather pattern that benefits from
vendor-specific fused kernels.  On CUDA, libraries such as Megablocks or
cuBLAS batched-GEMM with custom dispatch achieve near-zero overhead.
On ROCm, the same functionality is available via Composable Kernel (CK) and
hipBLASLt grouped GEMM, but the integration with PyTorch's dispatcher is
less mature.  For MiniMind's small ($N{=}4$) expert count and 768-dim hidden
size, the overhead is negligible, but larger MoE models would expose this
gap more sharply.

\paragraph{Serving ecosystem.}
vLLM and SGLang are CUDA-first; the ROCm ports exist but lag on features
such as FP8 KV-cache quantization, speculative decoding, and disaggregated
prefill-decode.  For production serving, these gaps translate to a real
efficiency penalty on AMD hardware that is separate from raw compute throughput.

\paragraph{Tooling and profiling.}
NVIDIA Nsight Systems/Compute provides fine-grained kernel-level profiling
that is tightly integrated with the CUDA runtime.  AMD's equivalent
(Omniperf, ROCprofiler-SDK, RocTracer) has reached parity for most common
workloads but has a steeper learning curve and less community-written tooling.

\subsection{Implications for the Research Community}
\label{sec:discussion:implications}

Our results provide three concrete takeaways:

\begin{enumerate}
\item \textbf{Algorithm development can and should be vendor-neutral.}
  The overhead of maintaining a portable device layer
  (\texttt{device\_utils.py}: 120 lines) is small relative to the algorithmic
  work in each trainer.  New algorithm proposals should publish on
  vendor-neutral code paths to avoid creating AMD-inaccessible baselines.

\item \textbf{The inference gap is the primary remaining concern.}
  Training throughput parity (where it exists) does not imply serving parity.
  The vLLM and SGLang ROCm feature gaps mean that AMD hardware may train
  efficiently but serve at reduced throughput until the inference ecosystem
  catches up.

\item \textbf{Small models are ideal for cross-vendor evaluation.}
  MiniMind's 64M parameters fit easily in a single-GPU context, which
  eliminates FSDP/TP parallelism as a confound.  Researchers validating
  vendor portability of new algorithms should prototype on small models first
  before investing in multi-node, mixed-precision, or expert-parallel stacks.
\end{enumerate}
